\documentclass[a4paper,12pt]{ctexart}
\usepackage[left=2.5cm, right=2.5cm, top=2.5cm, bottom=2.5cm]{geometry}
\usepackage{amsmath}
\usepackage{amssymb}
\usepackage{booktabs}
\usepackage{array}
\usepackage{graphicx}
\usepackage{float}
\usepackage{enumitem}
\usepackage{fancyhdr}

% --- 排版设置 ---
\linespread{1.5} % 1.5倍行距，方便书写
\setlength{\parskip}{0.5em}

% --- 定义题目环境 ---
\newcounter{probcount}
\newcommand{\problem}[1]{
    \par\noindent\textbf{题目 #1} \quad
}

% --- 定义书写空白命令 ---
% 默认留白 5cm，可根据需要调整
\newcommand{\workspace}[1][5cm]{
    \par
    \vspace{#1}
    \noindent\rule{\textwidth}{0.1pt} % 分隔线（可选，如果不需要可注释掉）
    \vspace{0.5em}
    \par
}

% --- 页眉页脚 ---
\pagestyle{fancy}
\fancyhf{}
\fancyhead[C]{数值计算方法 —— 专项刷题本}
\fancyfoot[C]{\thepage}

\title{\textbf{数值计算方法\\习题与例题汇编}}
\author{刷题本}
\date{\today}

\begin{document}

\maketitle

% 生成目录
\tableofcontents
\newpage

% ==================== 第一章 ====================
\section{第一章：误差与有效数字}

\subsection*{典型例题}

\problem{[3]}
要使 $\sqrt{20}$ 的近似值的相对误差限小于 $0.1\%$，要取几位有效数字？
\workspace

\problem{[4]}
已测得某场地长 $l$ 的值为 $l^* = 110\,\mathrm{m}$，宽 $d$ 的值为 $d^* = 80\,\mathrm{m}$，已知 $|l - l^*| \leqslant 0.2\,\mathrm{m}, |d - d^*| \leqslant 0.1\,\mathrm{m}$。试求面积 $s = ld$ 的绝对误差限与相对误差限。
\workspace

\subsection*{课后习题}

\problem{1}
设 $x>0$, $x$ 的相对误差为 $\delta$, 求 $\ln x$ 的误差.
\workspace

\problem{2}
设 $x$ 的相对误差为 $2\%$, 求 $x^n$ 的相对误差.
\workspace[4cm]

\problem{5}
计算球体体积要使相对误差限为 $1\%$, 问度量半径 $R$ 时允许的相对误差限是多少?
\workspace[4cm]

\problem{9}
正方形的边长大约为 100 cm, 应怎样测量才能使其面积误差不超过 1 $\text{cm}^2$?
\workspace

\newpage

% ==================== 第二章 ====================
\section{第二章：插值法}

\subsection*{典型例题}

\problem{[2]}
已给 $\sin 0.32 = 0.314567, \sin 0.34 = 0.333487, \sin 0.36 = 0.352274$，用线性插值及抛物插值计算 $\sin 0.3367$ 的值并估计截断误差。
\workspace[8cm]

\problem{[4]}
给出 $f(x)$ 的函数表（见下表），求 4 次牛顿插值多项式，并由此计算 $f(0.596)$ 的近似值。
\begin{table}[H]
\centering
\begin{tabular}{c c}
\toprule
$x$ & $f(x)$ \\
\midrule
$0.40$ & $0.41075$ \\
$0.55$ & $0.57815$ \\
$0.65$ & $0.69675$ \\
$0.80$ & $0.88811$ \\
$0.90$ & $1.02652$ \\
$1.05$ & $1.25382$ \\
\bottomrule
\end{tabular}
\end{table}
\workspace[8cm]

\subsection*{课后习题}

\problem{1}
当 $x = 1, -1, 2$ 时, $f(x) = 0, -3, 4$, 求 $f(x)$ 的二次插值多项式.
\begin{enumerate}
    \item 用单项式基底.
    \item 用拉格朗日插值基底.
    \item 用牛顿基底.
\end{enumerate}
证明三种方法得到的多项式是相同的.
\workspace[10cm]

\problem{2}
给出 $f(x) = \ln x$ 的数值表:
\begin{center}
\begin{tabular}{c|c|c|c|c|c}
\toprule
$x$ & 0.4 & 0.5 & 0.6 & 0.7 & 0.8 \\
\midrule
$\ln x$ & $-0.916\,291$ & $-0.693\,147$ & $-0.510\,826$ & $-0.356\,675$ & $-0.223\,144$ \\
\bottomrule
\end{tabular}
\end{center}
用线性插值及二次插值计算 $\ln 0.54$ 的近似值.
\workspace[8cm]

\problem{8}
$f(x) = x^7 + x^4 + 3x + 1$, 求 $f[2^0, 2^1, \cdots, 2^7]$ 及 $f[2^0, 2^1, \cdots, 2^8]$.
\workspace[4cm]

\newpage

% ==================== 第三章 ====================
\section{第三章：曲线拟合与逼近}

\subsection*{典型例题}

\problem{[9]}
已知一组实验数据如下表，求它的拟合曲线。
\begin{table}[H]
\centering
\begin{tabular}{c c c c c c}
\toprule
$x_i$ & $1$ & $2$ & $3$ & $4$ & $5$ \\
\midrule
$f_i$ & $4$ & $4.5$ & $6$ & $8$ & $8.5$ \\
$\omega_i$ & $2$ & $1$ & $3$ & $1$ & $1$ \\
\bottomrule
\end{tabular}
\end{table}
\workspace[8cm]

\problem{[10]}
设数据 $(x_i, y_i) \, (i=0, 1, 2, 3, 4)$ 由下表给出，可以看出数学模型为 $y = ae^{bx}$，用最小二乘法确定 $a$ 及 $b$。
\begin{table}[H]
\centering
\begin{tabular}{c c c c c c}
\toprule
$i$ & $0$ & $1$ & $2$ & $3$ & $4$ \\
\midrule
$x_i$ & $1.00$ & $1.25$ & $1.50$ & $1.75$ & $2.00$ \\
$y_i$ & $5.10$ & $5.79$ & $6.53$ & $7.45$ & $8.46$ \\
\bottomrule
\end{tabular}
\end{table}
\workspace[8cm]

\subsection*{课后习题}

\problem{4}
计算下列函数 $f(x)$ 关于 $C[0,1]$ 的 $\|f\|_\infty, \|f\|_1$ 与 $\|f\|_2$:
\begin{enumerate}
    \item $f(x) = (x-1)^3$;
    \item $f(x) = \left| x - \frac{1}{2} \right|$;
    \item $f(x) = x^m(1-x)^n, m$ 与 $n$ 为正整数.
\end{enumerate}
\workspace[10cm]

\problem{16}
观测物体的直线运动, 得出以下数据:
\begin{center}
\begin{tabular}{c|c|c|c|c|c|c}
\toprule
时间 $t/\text{s}$ & 0 & 0.9 & 1.9 & 3.0 & 3.9 & 5.0 \\
\midrule
距离 $s/\text{m}$ & 0 & 10 & 30 & 50 & 80 & 110 \\
\bottomrule
\end{tabular}
\end{center}
求运动方程.
\workspace[6cm]

\problem{17}
已知实验数据如下:
\begin{center}
\begin{tabular}{c|c|c|c|c|c}
\toprule
$x_i$ & 19 & 25 & 31 & 38 & 44 \\
\midrule
$y_i$ & 19.0 & 32.3 & 49.0 & 73.3 & 97.8 \\
\bottomrule
\end{tabular}
\end{center}
用最小二乘法求形如 $y = a + bx^2$ 的经验公式, 并计算均方误差.
\workspace[8cm]

\newpage

% ==================== 第四章 ====================
\section{第四章：数值积分与微分}

\subsection*{典型例题}

\problem{[1]}
给定形如 $\int_0^1 f(x)dx \approx A_0 f(0) + A_1 f(1) + B_0 f'(0)$ 的求积公式，试确定系数 $A_0, A_1, B_0$，使公式具有尽可能高的代数精度。
\workspace[6cm]

\problem{[3]}
对于函数 $f(x) = \frac{\sin x}{x}$，给出 $n=8$ 时的函数表（见下表），试用复合梯形公式及复合辛普森公式计算积分
$I = \int_0^1 \frac{\sin x}{x} dx$, 并估计误差。
\begin{table}[H]
\centering
\begin{tabular}{c c | c c}
\toprule
$x$ & $f(x)$ & $x$ & $f(x)$ \\
\midrule
$0$ & $1$ & $5/8$ & $0.9361556$ \\
$1/8$ & $0.9973978$ & $3/4$ & $0.9088516$ \\
$1/4$ & $0.9896158$ & $7/8$ & $0.8771925$ \\
$3/8$ & $0.9767267$ & $1$ & $0.8414709$ \\
$1/2$ & $0.9588510$ & & \\
\bottomrule
\end{tabular}
\end{table}
\workspace[8cm]

\subsection*{课后习题}

\problem{1}
确定下列求积公式中的待定参数, 使其代数精度尽量高, 并指明所构造出的求积公式具有的代数精度:
\begin{enumerate}
    \item $\int_{-h}^h f(x)dx \approx A_{-1}f(-h) + A_0f(0) + A_1f(h)$;
    \item $\int_{-2h}^{2h} f(x)dx \approx A_{-1}f(-h) + A_0f(0) + A_1f(h)$;
    \item $\int_{-1}^1 f(x)dx \approx [f(-1) + 2f(x_1) + 3f(x_2)]/3$;
    \item $\int_0^h f(x)dx \approx h[f(0) + f(h)]/2 + ah^2[f'(0) - f'(h)]$.
\end{enumerate}
\workspace[12cm]

\problem{4}
用辛普森公式求积分 $\int_0^1 e^{-x} dx$ 并估计误差.
\workspace[6cm]

\newpage

% ==================== 第五章 ====================
\section{第五章：线性方程组的直接解法}

\subsection*{典型例题}

\problem{[1]}
求 $A = \begin{pmatrix} 1 & -2 & 2 \\ -2 & -2 & 4 \\ 2 & 4 & -2 \end{pmatrix}$ 的特征值及谱半径。
\workspace[6cm]

\problem{[5]}
用直接三角分解法解
$\begin{pmatrix} 1 & 2 & 3 \\ 2 & 5 & 2 \\ 3 & 1 & 5 \end{pmatrix} \begin{pmatrix} x_1 \\ x_2 \\ x_3 \end{pmatrix} = \begin{pmatrix} 14 \\ 18 \\ 20 \end{pmatrix}$.
\workspace[8cm]

\subsection*{课后习题}

\problem{8}
用直接三角分解 (杜利特尔 Doolittle 分解) 求解线性方程组:
\[
\begin{cases}
\frac{1}{4}x_1 + \frac{1}{5}x_2 + \frac{1}{6}x_3 = 9, \\
\frac{1}{3}x_1 + \frac{1}{4}x_2 + \frac{1}{5}x_3 = 8, \\
\frac{1}{2}x_1 + x_2 + 2x_3 = 8
\end{cases}
\]
\workspace[8cm]

\newpage

% ==================== 第六章 ====================
\section{第六章：线性方程组的迭代解法}

\subsection*{典型例题}

\problem{[8]}
在线性方程组 $Ax = b$ 中，
$A = \begin{pmatrix} 1 & a & a \\ a & 1 & a \\ a & a & 1 \end{pmatrix}$,
证明当 $-\frac{1}{2} < a < 1$ 时高斯-塞德尔法收敛，而雅可比迭代法只在 $-\frac{1}{2} < a < \frac{1}{2}$ 时才收敛。
\workspace[8cm]

\subsection*{课后习题}

\problem{1}
设线性方程组
\[
\begin{cases}
5x_1 + 2x_2 + x_3 = -12, \\
-x_1 + 4x_2 + 2x_3 = 20, \\
2x_1 - 3x_2 + 10x_3 = 3.
\end{cases}
\]
(1) 考察用雅可比迭代法, 高斯-塞德尔迭代法解此方程组的收敛性; \\
(2) 用雅可比迭代法及高斯-塞德尔迭代法解此方程组, 要求当 $\|x^{(k+1)} - x^{(k)}\|_\infty < 10^{-4}$ 时迭代终止.
\workspace[10cm]

\problem{2}
设线性方程组
\[
(1) \begin{cases}
x_1 + 0.4x_2 + 0.4x_3 = 1, \\
0.4x_1 + x_2 + 0.8x_3 = 2, \\
0.4x_1 + 0.8x_2 + x_3 = 3;
\end{cases}
\quad
(2) \begin{cases}
x_1 + 2x_2 - 2x_3 = 1, \\
x_1 + x_2 + x_3 = 1, \\
2x_1 + 2x_2 + x_3 = 1.
\end{cases}
\]
试考察解此线性方程组的雅可比迭代法及高斯-塞德尔迭代法的收敛性.
\workspace[8cm]

\problem{4}
设 $A = \begin{bmatrix} 10 & a & 0 \\ b & 10 & b \\ 0 & a & 5 \end{bmatrix}$, $\det A \neq 0$, 用 $a, b$ 表示解线性方程组 $Ax=f$ 的雅可比迭代与高斯-塞德尔迭代收敛的充分必要条件.
\workspace[8cm]

\problem{6}
用雅可比迭代与高斯-塞德尔迭代解线性方程组
$Ax = b$, 证明若取 $A = \begin{bmatrix} 3 & 0 & -2 \\ 0 & 2 & 1 \\ -2 & 1 & 2 \end{bmatrix}$, 则两种方法均收敛, 试比较哪种方法收敛快?
\workspace[8cm]

\newpage

% ==================== 第七章 ====================
\section{第七章：非线性方程求根}

\subsection*{典型例题}

\problem{[2]}
求方程 $f(x) = x^3 - x - 1 = 0$ 在区间 $[1.0, 1.5]$ 内的一个实根，要求准确到小数点后的第 2 位。
\workspace[8cm]

\problem{[4]}
用不同方法求方程 $x^2 - 3 = 0$ 的根 $x^* = \sqrt{3}$。
\workspace[6cm]

\problem{[7]}
用牛顿法解方程 $xe^x - 1 = 0$.
\workspace[6cm]

\subsection*{课后习题}

\problem{1}
用二分法求方程 $x^3 - x - 1 = 0$ 的正根, 要求误差小于 $0.05$.
\workspace[8cm]

\problem{2}
为求方程 $x^3 - x^2 - 1 = 0$ 在 $x_0 = 1.5$ 附近的一个根, 设将方程改写成下列等价形式, 并建立相应的迭代公式.
(1) $x = 1 + 1/x^2$; \quad
(2) $x^3 = x^2 + 1$; \quad
(3) $x^2 = \frac{1}{x-1}$.
试分析每种迭代公式的收敛性, 并选出一种公式求出具有四位有效数字的近似根.
\workspace[10cm]

\newpage

% ==================== 第九章 ====================
\section{第九章：常微分方程数值解法}

\subsection*{典型例题}

\problem{[1]}
求解初值问题 $\begin{cases} y' = y - \frac{2x}{y}, \quad 0 < x < 1, \\ y(0) = 1. \end{cases}$
\workspace[8cm]

\subsection*{课后习题}

\problem{1}
用欧拉法解初值问题
\[
y' = x^2 + 100y^2, \quad y(0) = 0.
\]
取步长 $h=0.1$, 计算到 $x=0.3$ (保留到小数点后 4 位).
\workspace[8cm]

\end{document}